
\section{Detailed theoretical analysis}
\label{app:theory}

\subsection{Ordinal-probit parameterization}
\label{app:ordinal-probit}
Write \(\Delta(\mathcal R)\) for the probability simplex on
\(\mathcal R=\{0,1,2,3,4\}\).
The frozen Moment RM emits a latent location $m_\omega(x,y)$ and positive
scale $s_\omega(x,y)$.  With frozen cutpoints
$c_1<\cdots<c_4$, its five-category probabilities are
\(p_{\omega,0}=\Phi((c_1-m_\omega)/s_\omega)\),
\(p_{\omega,r}=\Phi((c_{r+1}-m_\omega)/s_\omega)-
\Phi((c_r-m_\omega)/s_\omega)\) for $r=1,2,3$, and
\(p_{\omega,4}=1-\Phi((c_4-m_\omega)/s_\omega)\).  The moments used by
the policy are \(\mu_\omega=\sum_{r=0}^4rp_{\omega,r}\) and
\(\sigma_\omega^2=\sum_{r=0}^4(r-\mu_\omega)^2p_{\omega,r}\).
Repeated-rating likelihood and disagreement supervise the conditional scale;
therefore this variance is a model-based signal of annotation disagreement,
not epistemic uncertainty or a stand-alone reward.  Proper scoring and
calibration are required because the complete categorical distribution enters
the finite-budget objective \citep{mccullagh1980regression,kendall2017uncertainties,
lou2024uncertainty,sun2025purm,gneiting2007strictly,guo2017calibration}.

This parameterization is included to make the evaluator interface reproducible,
not to claim a new ordinal-link model.  The methodological choice is what the
policy does with the frozen output: the same five probabilities support the
nominal mean, threshold products, robust envelope, and disagreement diagnostic.
Keeping these uses separate is important because replacing the vector by its
mean would make the finite-budget objective and its calibration check
unidentifiable.

\begin{proposition}[Finite-support threshold decomposition]
\label{prop:threshold-decomposition}
For conditionally independent ratings supported on
\(\mathcal R=\{0,1,2,3,4\}\),
the identity
\(h_N^{\rm nom}=\sum_{t=1}^{4}[1-\prod_{j=1}^{N}F_{\omega,j}(t-1)]\)
is exact, lies in \([0,4]\), and is nondecreasing under first-order stochastic
improvement of any candidate.  Its highest-threshold component is
\(1-\prod_j(1-p_{\omega,4}(x,Y_j))\).
\end{proposition}

The proposition is the finite-support sanity check for the main construction.
It also explains why the top-rating diagnostic in the artifact is not a
substitute for the reported objective: it is only the \(t=4\) term, whereas the
method trains against all four threshold events.  This distinction prevents a
policy from appearing successful by concentrating probability only at the
highest category while degrading the intermediate ordinal structure.

\begin{proof}
For any integer-valued \(Z\in\{0,\ldots,4\}\), the tail-sum identity gives
\(\mathbb E[Z]=\sum_{t=1}^{4}\Pr\{Z\ge t\}\).  Taking
\(Z=\max_jR_j\) and using (A1) gives
\(\Pr\{\max_jR_j<t\mid x,Y_{1:N}\}=\prod_{j=1}^{N}F_{\omega,j}(t-1)\).

Substitution proves the stated nominal identity.  Each term lies in
\([0,1]\), so the sum lies in \([0,4]\).  First-order
stochastic improvement decreases a candidate CDF and therefore cannot
decrease any summand.  At \(t=4\),
\(F_{\omega,j}(3)=1-p_{\omega,4}(x,Y_j)\), giving the highest-threshold
expression above.
\end{proof}

\begin{proposition}[Prompt-cluster calibration statement]
\label{prop:calibration-radius}
Assume that the validation prompt clusters are independent and that the
reward model is fixed before they are evaluated.  For cluster \(g\), let
\(R_{g\ell}\) denote its repeated ratings, \(n_g\) its number of ratings, and
\(F_{\omega,g\ell}(t)\) the frozen evaluator CDF at threshold \(t\).  Define
\(D_{g,t}=n_g^{-1}\sum_{\ell=1}^{n_g}
\{\mathbf{1}[R_{g\ell}\le t]-F_{\omega,g\ell}(t)\}\) for
\(t=0,\ldots,3\), and let
\(\widehat d_t=G_{\mathrm{cal}}^{-1}\sum_{g=1}^{G_{\mathrm{cal}}}D_{g,t}\), where
\(G_{\mathrm{cal}}\) is the number of independent validation clusters.  Let
\(\mathbb E_g\) denote expectation over the validation-cluster distribution.
Define the calibration radius as
\begin{equation}
\label{eq:epsilon-calibration}
\varepsilon_{\rm cal}
=
\min\!\left\{
1,\,
\max_{t=0,\ldots,3}(\widehat d_t)_+
+
\sqrt{
\frac{\log(4/\alpha_{\rm cal})}
     {2G_{\rm cal}}
}
\right\}.
\end{equation}
Then, with probability at least \(1-\alpha_{\rm cal}\),
\[
\mathbb E_g[D_{g,t}]\le\varepsilon_{\rm cal}
\qquad
\text{simultaneously for }t=0,1,2,3.
\]
\end{proposition}

\begin{proof}
Conditional on the frozen responses and predictions, changing all binary
indicators in a prompt cluster from zero to one changes \(D_{g,t}\) by
at most one.  Thus each cluster residual lies in an interval of width one.
Hoeffding's inequality gives
\(\Pr(\mathbb E_g[D_{g,t}]-\widehat d_t>
\sqrt{\log(4/\alpha_{\mathrm{cal}})/(2G_{\mathrm{cal}})})
\le\alpha_{\mathrm{cal}}/4\).

A union bound over four thresholds proves the simultaneous event.  On that
event, every right-hand side is bounded by the unclipped expression in
\cref{eq:epsilon-calibration}.  If clipping is active,
\(\varepsilon_{\mathrm{cal}}=1\), while
\(\mathbb E_g[D_{g,t}]\le1\), so the claim still holds.
\end{proof}

\emph{Scope of the calibration proposition.}
\Cref{prop:calibration-radius} is a population-average statement over the
validation prompt distribution, not a distribution-free pointwise guarantee
for every generated response.  The calibration split fixes the radius before
training, whereas the main theorem is conditional on its envelope covering a
particular candidate; the artifact records these two claims separately.

\begin{proof}[Proof of Theorem~\ref{thm:main-max}]
The envelope \(\overline F_j\) is nondecreasing, has boundary values zero and one, and therefore defines a valid categorical distribution. For every feasible \(q_j \in \mathcal U\), we have
\[
F_{q_j}(t) \le \overline F_j(t), \qquad t = 0,1,2,3.
\]
Hence \(q_j\) first-order stochastically dominates \(\overline p_j\), or equivalently \(\overline p_j\) is the stochastically smallest distribution in the ambiguity set. Since the maximum is coordinatewise nondecreasing in its ratings, replacing every feasible candidate distribution by \(\overline p_j\) cannot increase its expectation. The envelope itself is feasible and thus attains the infimum. Applying \Cref{prop:threshold-decomposition} gives the stated closed form.
\end{proof}

\begin{proof}[Proof of Theorem~\ref{thm:main-gradient}]
For nonnegative integer ratings, the positive-part marginal credit can be decomposed as
\[
(\overline R_i-\overline M_{-i})_+ = \sum_{t=1}^{4} \mathbf{1}\{\overline R_i \ge t,\; \overline M_{-i} < t\}.
\]
Taking conditional expectations and using the conditional independence assumption (A1) gives
\[
\Pr\{\overline R_i \ge t,\; \overline M_{-i} < t \mid x,Y_{1:N}\}
=
[1-\overline F_i(t-1)] \prod_{j \ne i} \overline F_j(t-1),
\]
which proves \Cref{eq:robust-credit-closed}.

For the gradient identity, fix \(i\) and condition on \((x, Y_{-i})\). The robust group value decomposes as
\[
h_N^{\mathrm{rob}} = \mathbb E[\overline M_{-i} \mid x, Y_{-i}] + q_i^{\mathrm{rob}}.
\]
The first term is independent of \(Y_i\), and the score-function identity implies that its product with \(\nabla_\theta \log \pi_\theta(Y_i \mid x)\) has conditional expectation zero. Applying the likelihood-ratio derivative to all \(N\) i.i.d. responses and summing the remaining terms proves \Cref{eq:robust-policy-gradient}.
\end{proof}

\begin{proof}[Proof of Theorem~\ref{thm:coverage-summary}]
Under (A1), let \(F_j^*\) denote the true conditional CDF. Under the pointwise
coverage premise \(F_j^*(t) \le \overline F_j(t)\) for every candidate and
threshold, first-order stochastic dominance directly yields the lower-bound
inequality stated in Theorem~\ref{thm:coverage-summary}. The same independence
assumption permits the product comparison used for the misspecified case.

For the misspecified case, define
\[
\xi \triangleq \max_{j,t}\, [F_j^*(t) - \overline F_j(t)]_+.
\]
At each threshold, the true and robust tail-sum terms can differ only through the corresponding CDF products. For arbitrary \(a_j, b_j \in [0,1]\), the product difference is bounded by the sum of positive differences:
\[
\prod_j a_j - \prod_j b_j \le \sum_j (a_j - b_j)_+.
\]
Applying this inequality with \(a_j = F_j^*(t-1)\) and \(b_j = \overline F_j(t-1)\) at each of the four thresholds, the discrepancy is at most \(N\xi\) per threshold, and therefore at most \(4N\xi\) after summation. Taking outer expectations preserves the resulting inequality:
\[
\mathbb E_*[\max_j R_j] \ge h_N^{\mathrm{rob}} - 4N\xi.
\]
\end{proof}

\begin{proof}[Proof of Proposition~\ref{prop:radius-monotonicity}]
Fix a candidate group and threshold \(t\), and define
\[
a_j(\varepsilon) \triangleq \min\{1,\; F_{\omega,j}(t-1) + \varepsilon\}.
\]
For \(0 \le \varepsilon_1 \le \varepsilon_2 \le 1\), the clipping map is nondecreasing and 1-Lipschitz, which gives
\[
0 \le a_j(\varepsilon_2) - a_j(\varepsilon_1) \le \varepsilon_2 - \varepsilon_1.
\]
Using the telescoping product identity and the fact that \(a_j(\varepsilon) \in [0,1]\), we have
\[
0 \le \prod_j a_j(\varepsilon_2) - \prod_j a_j(\varepsilon_1)
\le \sum_j \bigl(a_j(\varepsilon_2) - a_j(\varepsilon_1)\bigr)
\le N(\varepsilon_2 - \varepsilon_1).
\]
Substituting this bound into the four-term tail-sum expression for \(h_N^{\mathrm{rob}}\), and then taking expectations over prompts and responses, proves Proposition~\ref{prop:radius-monotonicity}.

At \(\varepsilon = 0\), the robust envelope equals the nominal CDF, so the exact infimum in Theorem~\ref{thm:main-max} equals \(J_N^{\mathrm{nom}}(\theta)\).
\end{proof}

\begin{proof}[Proof of Proposition~\ref{prop:variance}]
By the conditional credit definition in \Cref{eq:robust-credit-closed}, we have
\[
\mathbb E[(\overline R_i-\overline M_{-i})_+ \mid x, Y_{1:N}] = q_i^{\mathrm{rob}},
\]
for every \(i\). The score terms \(\phi_i\) are fixed under this conditioning, so linearity gives
\[
\mathbb E[G_{\mathrm{samp}} \mid x, Y_{1:N}] = G_{\mathrm{RB}}.
\]
For any fixed vector \(v\), the law of total variance yields
\[
\mathrm{Var}(v^\top G_{\mathrm{samp}})
=
\mathbb E\bigl[\mathrm{Var}(v^\top G_{\mathrm{samp}} \mid x, Y_{1:N})\bigr]
+
\mathrm{Var}(v^\top G_{\mathrm{RB}}),
\]
which directly implies the stated variance inequality.
\end{proof}

\begin{proof}[Proof of Proposition~\ref{prop:rollback}]
Let \((\theta_0, \mathsf{Opt}_0)\) denote the snapshot taken before a tentative update. If the measured rollout-state KL exceeds \(\delta_{\mathrm{hard}}\), the algorithm executes the restoration \((\theta, \mathsf{Opt}) \leftarrow (\theta_0, \mathsf{Opt}_0)\). By the semantics of the implementation's state-copy operation, both objects are restored exactly to their pre-update values. Because controller-state updates occur only in the accepted branch, a rejected transition also leaves the baseline, dual variables, KL coefficient, and RMS scale unchanged. This statement is conditional on the measured batch and therefore makes no population-level KL claim.
\end{proof}

\begin{proposition}[\(N=1\) reduction and complexity]
\label{cor:n1-reduction}
For \(N=1\),
\[
h_1^{\mathrm{rob}} = q_1^{\mathrm{rob}} = \sum_{r=0}^{4} r \overline p_{1,r}.
\]
Thus the upper-tail reward component reduces exactly to the robust mean reward. For general \(N\), all group values and marginal credits can be computed in \(O(N|\mathcal R|)\) time using threshold-wise prefix and suffix products.
\end{proposition}

This reduction serves as a useful implementation sanity check: at \(N=1\), the method must recover the robust mean and cannot manufacture a group advantage. For \(N>1\), the additional terms exactly capture the response's incremental chance of opening a new threshold above the other candidates. The test suite validates both regimes, ensuring that any improvement in the group objective cannot be attributed to a special-case implementation artifact.

\begin{proof}
With a single candidate, the empty opponent maximum is zero by convention, so \((\overline R_1 - 0)_+ = \overline R_1\), and the first equality follows directly. For the complexity analysis, note that evaluating each of the four nontrivial thresholds requires one forward and one backward product scan over the \(N\) candidates. Hence the total time is \(O(N \cdot |\mathcal R|)\), where \(|\mathcal R| = 5\).
\end{proof}
